% ---------- εξώφυλλο ----------
\thispagestyle{empty}
\begin{center}
\vspace*{26mm}
{\color{ggray}\small Η Γλώσσα των Μαθηματικών}\par\vspace{16mm}
{\bfseries\fontsize{24}{30}\selectfont Μαθηματικά Α’ Γυμνασίου\par}\vspace{9mm}
{\small Μέρος Α΄: Αριθμητική – Άλγεβρα \ · \ Μέρος Β΄: Γεωμετρία\par}\vspace{9mm}
{\color{gblue}\large Βοήθημα με νόημα: ασκήσεις που θέλεις να λύσεις\par}\vspace{40mm}
Ιωάννης Ψωμιάδης\par\vspace{2mm}
{\color{ggray}\scriptsize Σύμφωνα με το Σχολικό Βιβλίο της Α’ Γυμνασίου\par}\vspace{6mm}
{\small Οκτώβριος 2026\par}
\end{center}
\cleardoublepage
\pagestyle{fancy}
\fancyhead[RO]{\color{ggray}\small Πριν ξεκινήσεις}
\fancyhead[LE]{\color{ggray}\small Η Γλώσσα των Μαθηματικών}
\pagenumbering{roman}

\chapter*{}\vspace{-22mm}
{\color{gblue}\bfseries\Large Πριν ξεκινήσεις}\par\vspace{4mm}
{\color{gblue}\bfseries\large Ένα γράμμα σε σένα}\par

Μπήκες στο γυμνάσιο. Άλλο κτίριο, άλλα πρόσωπα, ένα πρόγραμμα με διαφορετικό μάθημα σε κάθε ώρα και έναν καθηγητή για καθένα. Αν έχεις διαβάσει τα βιβλία της Ε’ και της ΣΤ’, θα αναγνωρίσεις τα παιδιά, τον τρόπο και τη φιλοσοφία. Αν όχι, δεν πειράζει: όλα όσα χρειάζεσαι ξαναλέγονται εδώ, όπου χρειάζεται.

Τα περισσότερα βιβλία μαθηματικών σου λένε πώς λύνεται μια άσκηση. Αυτό το βιβλίο θέλει να σου πει κάτι άλλο: γιατί αξίζει να τη λύσεις.

Ξέρεις ήδη να μιλάς ελληνικά. Ίσως λίγο αγγλικά. Με αυτές τις γλώσσες περιγράφεις τον κόσμο: τι βλέπεις, τι θέλεις, τι σε ενοχλεί, τι σε κάνει να γελάς. Τα μαθηματικά είναι κι αυτά μια γλώσσα. Περιγράφουν τον ίδιο κόσμο, απλώς με άλλες λέξεις. Όταν γράφεις $5>3$, λες την ίδια πρόταση με το «το πέντε είναι μεγαλύτερο από το τρία». Το σύμβολο $>$ είναι μια λέξη. Το $=$ είναι μια λέξη. Ακόμα και ο αριθμός 1.000.000 είναι μια λέξη, που στα ελληνικά τη λέμε «ένα εκατομμύριο».

Στο δημοτικό έμαθες λέξεις. Στην ΣΤ’ το λεξιλόγιό σου μεγάλωσε. Φέτος μαθαίνεις γραμματική: τους κανόνες που λένε ποιες προτάσεις επιτρέπονται και ποιες όχι. Γιατί το $-3$ υπάρχει, αφού δεν μπορείς να έχεις «μείον τρία μήλα»; Γιατί $(-1)\cdot(-1)=1$; Γιατί ισχύει το «κάθε τρίγωνο έχει άθροισμα γωνιών 180°» και πώς το ξέρουμε χωρίς να τα μετρήσουμε όλα; Φέτος για πρώτη φορά δεν θα σε ρωτήσω μόνο «πόσο κάνει;». Θα σε ρωτήσω και «γιατί;».

Σε αυτό το βιβλίο θα μάθεις τη γλώσσα των μαθηματικών όπως μαθαίνει κανείς κάθε γλώσσα: λέξη λέξη, κανόνα κανόνα, με παραδείγματα από τη ζωή σου και με πολλή εξάσκηση. Δεν θα σου ζητήσω να αποστηθίσεις κανόνες που δεν καταλαβαίνεις. Κάθε κανόνας θα έχει έναν λόγο ύπαρξης, και θα τον δεις πριν δεις τον κανόνα.

Ένα ακόμη. Οι ασκήσεις εδώ δεν σου ζητούν να λύσεις κάτι επειδή έτσι πρέπει. Σου ζητούν να βοηθήσεις κάποιον, να κερδίσεις ένα στοίχημα, να αποδείξεις ότι έχεις δίκιο, να λύσεις μια διαφωνία. Τα μαθηματικά είναι το εργαλείο. Ο λόγος είναι πάντα δικός σου.

{\color{gblue}\bfseries\large Η παρέα}\par

Το Σεπτέμβρη η παρέα της ΣΤ’ βρέθηκε στο γυμνάσιο, στο ίδιο τμήμα, την Α3. Μαζί της βρέθηκαν και δύο παιδιά από το άλλο δημοτικό της περιοχής, και ένα κορίτσι που δεν ήξερε κανέναν. Οι δύο γειτονιές έγιναν μία τάξη. Η παρέα αυτή θα σε συνοδεύει σε όλο το γυμνάσιο, και τα προβλήματά της θα μεγαλώνουν μαζί της.

\begin{itemize}[leftmargin=1.4em,itemsep=0.2ex]
\item Η \textbf{Ζωή} ρωτάει πάντα «γιατί;» και δεν σταματάει στην πρώτη απάντηση.
\item Ο \textbf{Ηλίας} μετράει, δοκιμάζει, φτιάχνει πίνακες. Όταν κάνει λάθος, το βρίσκει μόνος του, αρκεί να τον αφήσεις.
\item Η \textbf{Μυρτώ} έχει το πρακτικό μυαλό της κουζίνας της γιαγιάς: αν μια απάντηση δεν βγάζει νόημα στην πραγματικότητα, δεν τη δέχεται.
\item Ο \textbf{Στέλιος} μαθαίνει έναν κανόνα και τον εφαρμόζει με ενθουσιασμό παντού, ακόμα και εκεί που δεν ταιριάζει.
\item Η \textbf{Φανή} κρατάει σημειώσεις για όλους και ρωτάει όσα οι άλλοι δεν τολμούν.
\item Ο \textbf{Βασίλης}, ο πρώην πρόεδρος του μαθητικού συμβουλίου, έχει πάντα άποψη και την υποστηρίζει με σιγουριά, κάποτε και με περισσότερη απ' όση χρειάζεται.
\item Ο \textbf{Νίκος}, από το άλλο δημοτικό, παίζει μπάσκετ και βιάζεται. Ό,τι γίνεται γρήγορα, το κάνει πρώτος. Ό,τι γίνεται σωστά, το κάνει δεύτερος.
\item Η \textbf{Αγγελική}, από το άλλο δημοτικό, ζωγραφίζει και παίζει μουσική. Βλέπει σχήματα και μοτίβα εκεί που οι άλλοι βλέπουν μόνο αριθμούς.
\item Η \textbf{Πολυξένη} είναι η καινούργια. Ήρθε το καλοκαίρι από τη Θεσσαλονίκη. Κρατάει ένα τετράδιο με κωδικούς και μοτίβα, μιλάει λίγο και ρωτάει σπάνια. Όταν ρωτάει, η ερώτησή της είναι αυτή που χρειαζόταν να ακουστεί.
\item Η \textbf{γιαγιά Ρήνη}, ο \textbf{παππούς} και ο \textbf{θείος Λευτέρης} είναι πάντα εκεί, με τον κήπο, το χωριό και τα πλακάκια τους.
\end{itemize}

Και υπάρχει η \textbf{κυρία Αλεξάνδρα}, η καθηγήτρια των μαθηματικών τους. Δεν δίνει ποτέ έτοιμη απάντηση. Δίνει ερώτηση.

\clearpage
{\color{gblue}\bfseries\large Πώς είναι φτιαγμένο το βιβλίο}\par

Όλες οι παράγραφοι έχουν την ίδια σειρά. Μόλις τη μάθεις, θα ξέρεις πάντα πού βρίσκεσαι.

\begin{enumerate}[leftmargin=2em,itemsep=0.5ex]
\item \textbf{Η ανάγκη.} Μια μικρή ιστορία με ένα πρόβλημα που δεν λύνεται χωρίς το εργαλείο του κεφαλαίου. Πρώτα νιώθεις ότι κάτι λείπει. Μετά σου το δίνω.
\item \textbf{Η λέξη.} Η νέα έννοια στα ελληνικά: τι σημαίνει, πού τη συναντάς. Μετά ο ακριβής ορισμός της.
\item \textbf{Η μετάφραση.} Ένας πίνακας με δύο στήλες. Αριστερά μια πρόταση στα ελληνικά, δεξιά η ίδια πρόταση στα μαθηματικά.
\item \textbf{Η θεωρία.} Οι κανόνες, με όλη την ακρίβεια που χρειάζονται. Είναι γραμμένοι σύμφωνα με το Σχολικό Βιβλίο σου, ώστε να βρίσκεις αμέσως την αντίστοιχη παράγραφο στην τάξη.
\item \textbf{Λυμένα παραδείγματα.} Κάθε λύση γράφεται ως σκέψη, σε τρία βήματα: Τι ξέρω. Τι θέλω. Πώς τα συνδέω. Στο τέλος ελέγχουμε αν η απάντηση βγάζει νόημα.
\item \textbf{Ασκήσεις σε τρεις ζώνες.}
  \begin{itemize}[leftmargin=1.2em,itemsep=0.2ex]
  \item \emph{Το κατάλαβα:} εφαρμόζεις τον κανόνα κατευθείαν.
  \item \emph{Το χρησιμοποιώ:} λύνεις ένα πρόβλημα από τη ζωή.
  \item \emph{Το σκέφτομαι:} μια ερώτηση που σε βάζει να σκεφτείς πιο μακριά, συχνά προς το επόμενο κεφάλαιο.
  \end{itemize}
\item \textbf{Λύσεις.} Όλες στο τέλος του βιβλίου, γραμμένες με τον ίδιο τρόπο σκέψης. Κοίτα τες αφού προσπαθήσεις, όχι πριν.
\end{enumerate}

Το βιβλίο χωρίζεται σε δύο μέρη, όπως και το Σχολικό σου: \textbf{Μέρος Α΄, Αριθμητική – Άλγεβρα}, και \textbf{Μέρος Β΄, Γεωμετρία}. Κάθε μέρος έχει κεφάλαια και κάθε κεφάλαιο \emph{παραγράφους} με την ίδια αρίθμηση που έχουν στο Σχολικό: η παράγραφος 1.4 του βιβλίου αυτού αντιστοιχεί στην Α.1.4 του Σχολικού. Οι ασκήσεις αριθμούνται με την παράγραφο: η άσκηση 1.4.7 είναι η έβδομη της παραγράφου 1.4.

Στο τέλος κάθε κεφαλαίου θα βρεις μια \textbf{Επανάληψη} με έναν «χάρτη»: ένα πίνακα που σου λέει σε ποιες παραγράφους να γυρίσεις, αν κάποια άσκηση σε δυσκόλεψε.

{\color{gblue}\bfseries\large Για τον γονέα}\par

Το βιβλίο αυτό γράφτηκε ώστε να το διαβάζετε μαζί με το παιδί σας. Δεν χρειάζεται να θυμάστε τα μαθηματικά του γυμνασίου.

Ξεκινήστε από την «Ανάγκη» και τη «Λέξη». Είναι γραμμένες σε απλά ελληνικά. Αν τις καταλάβετε εσείς, μπορείτε να τις εξηγήσετε. Η «Μετάφραση» δείχνει πώς η ίδια ιδέα γράφεται με σύμβολα, και εκεί συνήθως κολλάνε τα παιδιά: όχι στην ιδέα, αλλά στη μετάφρασή της.

Τρεις συμβουλές από δέκα χρόνια διδασκαλίας:

\begin{itemize}[leftmargin=1.4em,itemsep=0.4ex]
\item Ζητήστε από το παιδί να σας εξηγήσει μια λύση με δικά του λόγια. Αν μπορεί να την πει στα ελληνικά, την έχει καταλάβει.
\item Ένα λάθος είναι πληροφορία, όχι αποτυχία. Ρωτήστε «πώς το σκέφτηκες;» πριν πείτε «είναι λάθος».
\item Οι ασκήσεις της ζώνης «Το σκέφτομαι» δεν είναι υποχρεωτικές. Είναι για τις μέρες που υπάρχει όρεξη.
\end{itemize}

{\color{gblue}\bfseries\large Ποιος το έγραψε}\par

Σπούδασα μαθηματικά και δίδαξα ως φροντιστής για περίπου δέκα χρόνια. Τα επόμενα χρόνια δούλεψα στη βιομηχανία, κάνοντας μη καταστροφικούς ελέγχους: ελέγχω σωλήνες, δοχεία, δεξαμενές και συγκολλήσεις χωρίς να τα κόψουμε, για να ξέρουμε αν είναι ασφαλή.

Σε αυτά τα χρόνια έμαθα κάτι που δεν μου το είχε πει κανένα σχολικό βιβλίο. Τα μαθηματικά δεν είναι μάθημα. Είναι ο τρόπος που οι άνθρωποι συμφωνούν για το τι συμβαίνει. Όταν δύο μηχανικοί διαφωνούν, δεν κερδίζει αυτός που φωνάζει πιο δυνατά. Κερδίζει αυτός που φέρνει τη μέτρηση.

Αυτό το βιβλίο είναι η προσπάθειά μου να σου δώσω νωρίς αυτό το εργαλείο. Και, αν τα καταφέρω, λίγη από την αγάπη που έχω για αυτή τη γλώσσα.

\hfill Ιωάννης Ψωμιάδης

\cleardoublepage
\fancyhead[RO]{\color{ggray}\small Περιεχόμενα}
{\color{gblue}\bfseries\Large Περιεχόμενα}\par\vspace{3mm}
\makeatletter
\@starttoc{toc}
\makeatother
